\documentclass[11pt]{article}
\usepackage[ngerman]{babel}
\usepackage[a4paper,margin=0.9in]{geometry}
\usepackage[T1]{fontenc}
\usepackage[utf8]{inputenc}
\usepackage{lmodern}
\usepackage{amsmath,amssymb,amsthm,mathtools,mathrsfs}
\usepackage{microtype}
\usepackage{fancyhdr}
\usepackage{array,longtable,enumitem,multicol}
\usepackage{tikz}
\setlength{\parindent}{1.5em}
\setlength{\parskip}{0.18em}
\setlength{\headheight}{14pt}
\emergencystretch=8em
\pagestyle{fancy}
\fancyhf{}
\cfoot{\thepage}
\lhead{Heinrich Weber, Lehrbuch der Algebra, Vol. II}
\newcommand{\eps}{\varepsilon}
\newcommand{\ii}{\mathrm{i}}
\newcommand{\Om}{\Omega}
\newcommand{\tg}{\operatorname{tg}}
\newcommand{\PSL}{\operatorname{PSL}}
\newcommand{\Jac}{\operatorname{Jac}}
\newcommand{\mat}[1]{\begin{pmatrix}#1\end{pmatrix}}
\newcommand{\arr}[1]{\begin{array}{#1}}
\newcommand{\G}{G_{168}}
\rhead{German Source}
\begin{document}

\section*{\S\,121. Invariantencurven der Gruppe \(G_{168}\).}

Eine Form \(n\)-ten Grades
\begin{equation}
\varphi(x_1,x_2,x_3)
\end{equation}
wird durch die Substitutionen der Gruppe im Allgemeinen in 168 verschiedene Formen uebergehen. Bei besonderen Formen von \(\varphi\) kann diese Zahl sich aber verringern; dann bilden die Substitutionen, durch die \(\varphi\) ungeaendert bleibt, eine in \(G_{168}\) enthaltene Gruppe \(G'\), deren Grad ein Teiler von 168 sein muss. Die Anzahl der Formen, in die \(\varphi\) uebergeht, ist der Index des Teilers \(G'\), und jede dieser verschiedenen Formen von \(\varphi\) bleibt durch eine mit \(G'\) conjugirte Gruppe ungeaendert. Die Form \(\varphi\) ist eine absolute Invariante der Gruppe \(G'\).

Die Gleichung \(\varphi=0\) stellt eine auf das Coordinatendreieck \(x_1,x_2,x_3\) bezogene Curve dar, und diese Curve wird durch die Substitutionen von \(G_{168}\) auf andere Curven abgebildet. Sind die 168 Bildcurven nicht alle von einander verschieden, so bleibt die Curve \(\varphi\) durch die Substitutionen einer Gruppe \(G'\) ungeaendert, oder sie aendert sich nur um einen constanten Factor; sie ist also eine relative Invariante von \(G'\).

Eine gerade Linie bleibt nur dann durch eine von der identischen verschiedene Substitution ungeaendert, wenn sie zu den im vorigen Paragraphen beschriebenen Axen gehoert. Alle Eigenschaften und Beziehungen zwischen Punkten, Linien und Curven, die durch lineare Transformationen unzerstoerbar sind, also die projectiven Eigenschaften der Geometrie, bleiben in den Bildern erhalten. Ist also ein Punkt einer Curve ein Doppelpunkt, Wendepunkt, Rueckkehrpunkt, Beruehrungspunkt einer Doppeltangente u. s. w., so besitzen die Bildpunkte die entsprechende Eigenschaft.

Unter den Formen \(\varphi\) sind vor allem diejenigen wichtig, die bei der ganzen Gruppe ungeaendert bleiben: die Invarianten. Da die Gruppe einfach ist, gibt es hier nur absolute Invarianten. Ist \(\varphi(x_1,x_2,x_3)\) eine solche Form, so heisst die durch \(\varphi=0\) dargestellte Curve eine invariante Curve der Gruppe. Es gibt dann 168 Abbildungen der Ebene auf sich selbst, bei denen alle Punkte der Curve in Punkte derselben Curve uebergehen. Liegt irgend ein Punkt auf dieser Curve, so liegen auch alle seine Bildpunkte darauf.

Im Allgemeinen ist die Zahl der miteinander verbundenen Punkte einer solchen Curve 168, jedenfalls nicht groesser. Ist sie kleiner, so muss der Punkt ein Pol sein, und die Zahl der Bildpunkte ist ein Teiler von 168, naemlich 24 fuer die \(P_7\), 21 fuer die \(P_8\), 28 fuer die \(P_6\), 42 fuer die \(P_4\), 56 fuer die \(P_3\), und 84 fuer ein System zusammengehoeriger \(P_2\). Wenn ein Pol einer dieser Arten auf der Curve liegt, so liegen alle Pole derselben Art darauf.

\section*{\S\,122. Die erste Invariante der Gruppe \(G_{168}\) und die Grundcurve.}

Um die Invarianten unserer Gruppe zu bilden, suchen wir Formen der drei Variablen \(x_1,x_2,x_3\), die durch Anwendung der drei Substitutionen \(\chi,r,\omega\) ungeaendert bleiben. Dies genuegt, da die ganze Gruppe aus diesen drei Substitutionen zusammengesetzt wird. Die Substitution \(\chi\) ist eine cyklische Vertauschung der drei Variablen, und \(r\) bedeutet
\begin{equation}
\mat{x_1&x_2&x_3\\ \eps x_1&\eps^2x_2&\eps^4x_3}.
\end{equation}
Wenn also in einer Invariante ein Glied \(x_1^{h_1}x_2^{h_2}x_3^{h_3}\) vorkommt, worin \(h_1,h_2,h_3\) nicht negative ganze Zahlen sind, so verlangt die Unveraenderlichkeit durch \(r\) die Congruenz
\begin{equation}
h_1+2h_2+4h_3\equiv0\pmod 7,
\tag{1}
\end{equation}
und die Unveraenderlichkeit durch \(\chi\), dass neben diesem einen Gliede die zwei entsprechenden cyklischen Glieder vorkommen, wenn nicht \(h_1=h_2=h_3\) ist. Dazu kommt noch die Bedingung der Unveraenderlichkeit durch \(\omega\).

Wir bestimmen zuerst den kleinsten Grad \(m=h_1+h_2+h_3\), fuer den dies moeglich ist. Die Bedingung (1) kann fuer \(m<3\) nicht erfuellt sein. Fuer \(m=3\) muesste \(h_1=h_2=h_3=1\) sein; das Product \(x_1x_2x_3\) bleibt aber durch \(\omega\) nicht ungeaendert. Der naechste Fall ist \(m=4\). Aus
\begin{equation}
h_1+h_2+h_3=4,
\qquad
h_1+2h_2+4h_3\equiv0\pmod 7
\end{equation}
folgen, bis auf cyklische Vertauschung, die Exponenten
\begin{equation}
(3,0,1),\qquad (0,3,1),\qquad (1,0,3).
\end{equation}
Daher ergibt sich, von einem Factor abgesehen, die einzige durch \(r\) und \(\chi\) ungeaenderte Form vierten Grades
\begin{equation}
f(x_1,x_2,x_3)=x_1^3x_3+x_2^3x_1+x_3^3x_2.
\tag{2}
\end{equation}

Um den Einfluss von \(\omega\) auf \(f\) zu pruefen, setzen wir
\begin{equation}
\begin{aligned}
x_1&=\alpha y_1+\beta y_2+\gamma y_3,\\
x_2&=\beta y_1+\gamma y_2+\alpha y_3,\\
x_3&=\gamma y_1+\alpha y_2+\beta y_3,
\end{aligned}
\tag{3}
\end{equation}
und nehmen an, dass dadurch
\begin{equation}
F(y_1,y_2,y_3)=f(x_1,x_2,x_3)
\tag{4}
\end{equation}
geleistet werde. Bei der Bildung der zweiten Ableitungen benutzt man
\begin{equation}
\begin{array}{lll}
\frac16 f_{11}=x_1x_3,&\frac16 f_{22}=x_1x_2,&\frac16 f_{33}=x_2x_3,\\[2mm]
\frac13 f_{23}=x_3^2,&\frac13 f_{31}=x_1^2,&\frac13 f_{12}=x_2^2.
\end{array}
\end{equation}
Die aus (3) aufgeloesten Gleichungen haben dieselbe cyklische Form. Mit Hilfe der Relationen zwischen \(\alpha,\beta,\gamma\) aus \S\,115 verschwinden die zusaetzlichen Coefficienten, und man erhaelt wieder genau die entsprechenden Ableitungen von \(f(y_1,y_2,y_3)\). Nach dem Euler'schen Satze folgt
\begin{equation}
F(y_1,y_2,y_3)=y_1^3y_3+y_2^3y_1+y_3^3y_2.
\end{equation}
Damit ist gezeigt, dass \(f\) wirklich eine Invariante der Gruppe \(G_{168}\) ist. Sie ist, abgesehen von einem constanten Factor, die einzige Invariante vierten Grades.

Die Gleichung
\begin{equation}
f(x_1,x_2,x_3)=0
\tag{5}
\end{equation}
bedeutet eine Curve vierter Ordnung; wir nennen sie die Grundcurve der Gruppe oder die Curve \(f\). Es gibt 168 Abbildungen der Ebene auf sich selbst, bei denen jedem Punkte dieser Curve ein Punkt derselben Curve entspricht.

Die gerade Linie \(x_1=0\) schneidet die Curve in drei zusammenfallenden Punkten und in einem vierten davon getrennten Punkte. Die Eckpunkte des Coordinatendreiecks sind also Wendepunkte der Curve, und die Seiten sind die Wendetangenten. Bezeichnen wir die Seiten des Coordinatendreiecks mit \(1,2,3\) und die gegenueberliegenden Ecken durch dieselben Ziffern, so ist Seite \(1\) Wendetangente im Punkte \(2\), Seite \(2\) Wendetangente im Punkte \(3\), und Seite \(3\) Wendetangente im Punkte \(1\).

Wir untersuchen nun die Lage der Pole und Axen in Bezug auf die Grundcurve. Da die Eckpunkte des Coordinatendreiecks zu den siebenzaehligen Polen gehoeren, liegen alle \(P_7\) auf der Grundcurve. Sie sind, wie die Eckpunkte selbst, Wendepunkte der Curve und ordnen sich in acht Wendepunktsdreiecke.

Auch die dreizaehligen Pole \(P_3\) liegen auf der Grundcurve. Setzt man, mit einer cubischen Einheitswurzel \(\rho\),
\begin{equation}
x_2=\rho x_1,
\qquad
x_3=\rho^2x_1,
\qquad
1+\rho+\rho^2=0,
\end{equation}
so erfuellt der Punkt die Gleichung der Grundcurve. Die Verbindungslinie des zugehoerigen Paares ist
\begin{equation}
T=x_1+x_2+x_3=0.
\end{equation}
Setzt man in (2) \(x_1=-x_2-x_3\), so geht \(f=0\) in
\begin{equation}
\bigl(x_2^2+x_2x_3+x_3^2\bigr)^2=0
\end{equation}
ueber. Die Linie \(T\) ist also eine Doppeltangente. Die 28 Linien \(T\) sind Doppeltangenten der Grundcurve; ihre Beruehrungspunkte sind die 56 Pole \(P_3\). Die Pole \(P_7\) koennen nicht auf diesen Linien liegen.

Fuer die Punkte \(P_6\) ist die Nichtzugehoerigkeit zur Grundcurve unmittelbar zu sehen, wenn man \(x_1=x_2=x_3=1\) setzt; dann wird \(f=3\). Fuer die \(P_8\) erhaelt man aus der Darstellung in \S\,118 und den Relationen in \S\,115 ebenfalls einen von Null verschiedenen Wert von \(f\). Die Schnittpunkte der 21 Hauptaxen untereinander sind die Pole \(P_8\) und \(P_6\), daher schneiden sich niemals zwei Hauptaxen auf der Grundcurve. Die Schnittpunkte einer Hauptaxe mit der Grundcurve koennen also nur vierzaehlige oder zweizaehlige Pole sein. Wuerde ein \(P_4\) auf einer Hauptaxe und auf \(f\) liegen, so muesste die Axe eine Doppeltangente sein; da die 28 Doppeltangenten aber schon die Linien \(T\) sind, ist dies ausgeschlossen. Also sind die Schnittpunkte der Grundcurve mit den Hauptaxen zweizaehlige Pole.

Damit ergeben sich die ausgezeichneten Punktsysteme
\begin{equation}
24P_7,\\ 56P_3,\\ 84P_2\quad\text{auf der Curve }f,
\qquad
21P_8,\\42P_4,\\28P_6\quad\text{nicht auf der Curve }f.
\end{equation}
Alle anderen Punkte der Curve \(f\) gehen durch die Substitutionen der Gruppe in 168 verschiedene Punkte ueber. Daraus folgt auch, dass die Curve \(f\) keinen Doppelpunkt besitzt und nicht in Curven niedrigeren Grades zerfaellt. Haette sie naemlich einen Doppelpunkt, so muessten auch alle seine Bildpunkte Doppelpunkte sein; eine Curve vierter Ordnung kann aber, wenn sie keinen doppelt gezaehlten Curventheil besitzt, nicht mehr als sechs Doppelpunkte haben. Andernfalls waere \(f\) das Quadrat einer quadratischen Form, was wegen des Fehlens quadratischer Invarianten unmoeglich ist. Ein System von Punkten, deren jeder durch Substitutionen der Gruppe in jeden anderen uebergefuehrt werden kann, nennen wir ein System verbundener Punkte.

\section*{\S\,123. Die hoeheren Invarianten.}

Weitere Invarianten unserer Gruppe lassen sich nach dem Satze aus \S\,40 ableiten, indem man Covarianten der Form \(f\) bildet. Wir bilden zuerst die Hesse'sche Covariante
\begin{equation}
\Delta={1\over54}
\begin{vmatrix}
f_{11}&f_{12}&f_{13}\\
f_{21}&f_{22}&f_{23}\\
f_{31}&f_{32}&f_{33}
\end{vmatrix}.
\end{equation}
Fuer (2) erhaelt man, bis auf einen constanten Factor,
\begin{equation}
\Delta=x_1^5x_2+x_2^5x_3+x_3^5x_1-5x_1^2x_2^2x_3^2.
\tag{3}
\end{equation}
Auf der Curve \(\Delta=0\) liegen die Ecken des Coordinatendreiecks, folglich alle siebenzaehligen Pole \(P_7\). Sie bilden das vollstaendige System der Schnittpunkte von \(f\) und \(\Delta\).

Eine weitere Covariante, und zwar vom vierzehnten Grade, erhalten wir aus der Determinante
\begin{equation}
C={1\over9}
\begin{vmatrix}
f_{11}&f_{12}&f_{13}&\Delta_1\\
f_{21}&f_{22}&f_{23}&\Delta_2\\
f_{31}&f_{32}&f_{33}&\Delta_3\\
\Delta_1&\Delta_2&\Delta_3&0
\end{vmatrix}.
\end{equation}
Sie laesst sich aus (3) berechnen und hat, bei der vorliegenden Normierung, die Form
\begin{equation}
\begin{aligned}
C={}&x_1^{14}+x_2^{14}+x_3^{14}\\
&-34\bigl(x_1^{11}x_2x_3^2+x_1^2x_2^{11}x_3+x_1x_2^2x_3^{11}\bigr)\\
&-250\bigl(x_1^9x_2^4x_3+x_1^4x_2x_3^9+x_1x_2^9x_3^4\bigr)\\
&+375\bigl(x_1^8x_2^2x_3^4+x_1^4x_2^8x_3^2+x_1^2x_2^4x_3^8\bigr)\\
&+18\bigl(x_1^7x_2^7+x_2^7x_3^7+x_3^7x_1^7\bigr)\\
&-126\bigl(x_1^6x_2^5x_3^3+x_1^5x_2^3x_3^6+x_1^3x_2^6x_3^5\bigr).
\end{aligned}
\tag{4}
\end{equation}
Schon aus dem ersten Gliede erkennt man, dass die Curve \(C=0\) nicht durch die Ecken des Coordinatendreiecks geht.

Eine vierte Invariante, vom einundzwanzigsten Grade, erhaelt man als Functionaldeterminante der drei Formen \(f,\Delta,C\):
\begin{equation}
K={1\over14}
\begin{vmatrix}
f_1&\Delta_1&C_1\\
f_2&\Delta_2&C_2\\
f_3&\Delta_3&C_3
\end{vmatrix}.
\tag{5}
\end{equation}
Auch diese Curve geht nicht durch die Ecken des Coordinatendreiecks. Die vollstaendig ausgerechneten Ausdruecke fuer \(C\) und \(K\) sind die bekannten Gordan'schen Formen fuer die ternäre biquadratische Form \(f\).

\section*{\S\,124. Das volle Invariantensystem.}

Wir koennen nun nachweisen, dass die Formen \(f,\Delta,C,K\) ein volles Invariantensystem der Gruppe sind, d. h. dass alle Invarianten der Gruppe als ganze rationale Functionen dieser vier Formen dargestellt werden koennen. Wir stuetzen uns dabei auf den Satz von Bezout: zwei Curven der Ordnungen \(m\) und \(n\), die mehr als \(mn\) Punkte gemeinsam haben, besitzen einen gemeinsamen Curventheil; Beruehrungspunkte sind dabei doppelt oder mehrfach zu zaehlen.

Sei \(\Phi\) eine Invariante \(m\)-ter Ordnung, die \(f\) nicht als Factor enthaelt. Unter den Schnittpunkten der Curven \(\Phi=0\) und \(f=0\) moegen \(h_1\) mal die Pole \(P_7\), \(h_2\) mal die Pole \(P_3\), \(h_3\) mal die Pole \(P_2\) vorkommen; ausserdem moegen noch \(h_4\) Systeme von je 168 verbundenen Punkten auftreten. Dann ist
\begin{equation}
4m=24h_1+56h_2+84h_3+168h_4,
\end{equation}
und daher
\begin{equation}
m=6h_1+14h_2+21h_3+42h_4.
\tag{1}
\end{equation}
Die Zahlen \(h_i\) sind nicht negative ganze Zahlen und koennen nicht alle verschwinden.

Der kleinste moegliche Grad ist also \(m=6\). Eine Invariantencurve sechster Ordnung muss durch die 24 Punkte \(P_7\) gehen, wie dies fuer \(\Delta\) gilt. Eine zweite Invariante sechster Ordnung unterscheidet sich daher von \(\Delta\) nur um einen Ausdruck, der durch \(f\) teilbar waere; da keine Invariante zweiten Grades vorhanden ist, ist sie proportional zu \(\Delta\). Ebenso zeigen die Faelle der Grade 12 und 18, dass die betreffenden Invarianten rationale Functionen von \(\Delta\) und \(f\) sind.

Der naechste Fall ist \(m=14\). Dann muessen alle Invarianten vierzehnter Ordnung durch die 56 Punkte \(P_3\) gehen; diese bilden das vollstaendige Schnittpunktsystem einer solchen Curve mit der Grundcurve. Daher hat jede Invariante vierzehnter Ordnung die Form
\begin{equation}
aC+bf\Delta,
\end{equation}
und umgekehrt ist jeder Ausdruck dieser Form eine Invariante vierzehnter Ordnung.

In gleicher Weise behandelt man die naechsten Grade. Eine unabhaengige Invariante zwanzigster Ordnung tritt nicht auf. Fuer den Grad 21 gibt es, von einem constanten Factor abgesehen, nur eine Invariante. Das System der 21 Hauptaxen ist aber eine Invariante einundzwanzigster Ordnung; folglich gilt:
\begin{quote}
Die Invariante \(K\) zerfaellt in 21 lineare Factoren, die, gleich Null gesetzt, die Hauptaxen der Gruppe darstellen.
\end{quote}

Sodann bildet man eine Invariante zweiundvierzigsten Grades
\begin{equation}
Q=\Delta^7-kC^3,
\tag{2}
\end{equation}
wobei die Constante \(k\) passend gewaehlt wird. Sie kann so bestimmt werden, dass die Curve \(Q=0\) ein beliebig gegebenes System verbundener Punkte aus der Grundcurve ausschneidet. Ist nun \(\Phi\) eine beliebige durch \(f\) nicht teilbare Invariante, so bildet man zu den auf \(f\) liegenden verbundenen Schnittpunktsystemen die entsprechenden Formen \(Q_1,Q_2,\ldots\) und setzt
\begin{equation}
W=\Phi-a\Delta^{h_1}C^{h_2}K^{h_3}Q_1Q_2\cdots.
\end{equation}
Die Constante \(a\) kann so gewaehlt werden, dass \(W\) durch einen weiteren Punkt der Grundcurve geht; dann ist \(W\) nach Bezout durch \(f\) teilbar. Der Quotient ist wieder eine Invariante niedrigeren Grades. Vollstaendige Induction liefert daher den Satz:
\begin{quote}
Jede Invariante der Gruppe laesst sich als ganze rationale Function der vier fundamentalen Invarianten \(f,\Delta,C,K\) darstellen.
\end{quote}

Waehlt man in (2) die Constante so, dass \(Q\) durch einen der Pole \(P_2\) geht, so folgt weiter, dass \(K^2-k'Q\) durch \(f\) teilbar wird. Der Quotient ist eine Invariante geraden Grades und kann also ohne \(K\) durch die fundamentalen Invarianten dargestellt werden. Somit kann \(K^2\) rational durch \(f,\Delta,C\) ausgedrueckt werden. In der Form
\begin{equation}
K^2=\sum a_{\nu,\nu_1,\nu_2} f^\nu\Delta^{\nu_1}C^{\nu_2},
\qquad
2\nu+3\nu_1+7\nu_2=21,
\tag{3}
\end{equation}
besteht eine lineare Relation mit numerischen Coefficienten. Die Relation gestattet, in jedem Ausdruck alle Potenzen von \(K\) mit Ausnahme der ersten zu eliminiren. Daraus folgt endlich:
\begin{quote}
Eine Invariante geraden Grades laesst sich rational durch \(f,\Delta,C\) darstellen; eine Invariante ungeraden Grades ist das Product von \(K\) mit einer Invariante geraden Grades.
\end{quote}

\clearpage
\end{document}
